\begin{abstract}
A multi-task leaderboard ranks models by aggregating per-task scores, and developers can try to climb it by training on data related to some of the sub-datasets. We ask how much such targeted investment a developer needs to make one target model rank first, and how this depends on the aggregation rule. We model investment as raising the target's scores on chosen tasks within a cap, with all other models fixed, and measure it by the number of tasks touched ($K$), the total points raised ($L_1$) and the largest single-task jump ($L_\infty$), always under a named cap scenario and together with the rate at which first place is attainable. For score rules we derive closed forms (mean, minimum, odd-$n$ median) and solve power means and rank rules (Borda, plurality, Condorcet winner) exactly by mixed-integer programming. On public data (HELM MMLU v1.13.0, 79 models $\times$ 57 subjects; Open LLM Leaderboard, 4496 models $\times$ 6 tasks) the answer is conditional: for the ten closest challengers on HELM, the median-rule, minimum-rule and plurality costs are only 3--6\% of the mean-rule cost in total points (paired medians), Borda and the Condorcet requirement are 41\% and 17\%, and geometric and harmonic means are 91\% and 84\%; the ordering of rules by cost is nearly unchanged across six cap scenarios, whereas first-place feasibility is not (the minimum rule is attainable for only 6 and 7 of 10 targets under the two tightest gain caps). Under evaluation noise the cheap rules have the largest relative spread of cost, though in points their spread is small. Contrary to our expectation, order-statistic rules force little breadth for near-top targets: with paired targets the median rule's $K$ grows with the number of tasks (log--log slope of the mean cost 0.54, interval 0.51--0.56, against 0.84 for the mean rule) but stays far below the worst case of a majority of tasks, and the minimum rule's barely (0.16). On the six-task Open LLM Leaderboard the score scale matters, in paired comparisons over targets: the median rule's cost relative to the mean rule rises from 0.13 to 0.84 (ranks 11--30; paired difference 0.70, interval 0.61--0.82) and from 0.34 to 0.78 (ranks 31--100) when normalised scores are replaced by raw ones, but not beyond rank 100 (+0.008), and the minimum rule is dearer on raw scores in every rank group. We state where the conclusions fail and what we did not complete.
\end{abstract}
